\section{Kriptografi Simetris dan Asimetris}

Berdasarkan penggunaan kunci, kriptografi dibagi menjadi dua kategori utama:

\begin{table}[!htbp]
\centering
\caption{Perbandingan kriptografi simetris dan asimetris}
\label{tab:simetris-asimetris}
\begin{tabular}{@{}lll@{}}
\toprule
\textbf{Aspek} & \textbf{Simetris} & \textbf{Asimetris} \\
\midrule
Jumlah kunci & 1 kunci rahasia & 2 kunci (publik + privat) \\
Enkripsi/Dekripsi & Kunci sama & Kunci berbeda \\
Contoh & AES, DES, Caesar & RSA, ElGamal \\
Kecepatan & Cepat & Lebih lambat \\
Distribusi kunci & Sulit (perlu saluran aman) & Mudah (kunci publik terbuka) \\
\bottomrule
\end{tabular}
\end{table}

\begin{figure}[!htbp]
\centering
\begin{tikzpicture}[
  node distance=1cm,
  box/.style={rectangle, draw, minimum width=2cm, minimum height=0.7cm, align=center, font=\small}
]
\node[box, fill=blue!20] (alice) {Alice};
\node[box, fill=green!20, right=4cm of alice] (bob) {Bob};
\node[above=0.5cm of alice, font=\scriptsize] {Kunci K};
\node[above=0.5cm of bob, font=\scriptsize] {Kunci K};
\draw[-Stealth, thick] (alice) -- node[above] {ciphertext} (bob);
\draw[<->, dashed, blue] (alice.north) -- ++(0,0.8) -| (bob.north) node[midway, above, font=\scriptsize] {Kunci rahasia};
\end{tikzpicture}
\caption{Kriptografi simetris: satu kunci untuk enkripsi dan dekripsi}
\label{fig:simetris}
\end{figure}

\subsection{Kriptografi Simetris}
Menggunakan \textbf{satu kunci} yang sama untuk enkripsi dan dekripsi \cite{katz2020,trappe2006}. Alice dan Bob harus berbagi kunci rahasia secara aman sebelum berkomunikasi. Contoh: Caesar cipher, AES, DES. Kelebihan: cepat dan efisien. Kelemahan: masalah distribusi kunci---bagaimana mengirim kunci secara aman?

\subsection{Kriptografi Asimetris}
Menggunakan \textbf{sepasang kunci}: kunci publik (untuk enkripsi) dan kunci privat (untuk dekripsi). Kunci publik dapat didistribusikan secara terbuka; hanya kunci privat yang dijaga rahasia. Contoh: RSA. Kelebihan: tidak perlu saluran aman untuk pertukaran kunci. Kelemahan: lebih lambat dibanding kriptografi simetris \cite{schneier1996}.

\subsection{Kriptografi Hibrida}
Dalam praktik, sistem modern sering menggabungkan keduanya: kriptografi asimetris untuk pertukaran kunci, lalu kriptografi simetris untuk mengenkripsi data aktual (misalnya TLS/SSL).
